\section{Résultats attendus et calendrier}

Le projet se déroulera durant le trimestre d'hiver de l'année académique 2026--2027. Au-del\`a du mentorat individuel, plusieurs activités rassembleront l'ensemble des participant(e)s (mentors et mentoré(e)s) afin de favoriser les échanges autour des mathématiques et de renforcer la cohésion du groupe. Nos objectifs sont :

\begin{itemize}[leftmargin=1.6em,labelsep=0.6em]

    \item[\ding{118}] Créer un environnement qui stimule la participation proactive et soutenue des mentors et des étudiant(e)s aux projets, selon leurs disponibilités respectives. L'expérience de l'année dernière suggère qu'une rencontre hebdomadaire est une bonne solution, mais afin de répondre aux besoins de chacun(e), elle ne sera imposée à personne.

    \item[\ding{118}] Organiser des rencontres mensuelles réunissant l'ensemble des participant(e)s afin de créer une synergie au sein du groupe.

    \item[\ding{118}] Dans l'esprit de dialogue du projet, des rencontres seront organisées avec des professeur(e)s qui, endossant le rôle de mentors, exposeront certains aspects (méthodologiques, professionnels) ainsi que leurs impressions sur leur travail et sur son évolution au fil du temps (rôle de l'IA, alternatives à la recherche académique), au bénéfice tant des mentors que des étudiant(e)s.

    \item[\ding{118}] Chaque mentoré(e) devra rédiger un rapport d'au plus 15 pages et présenter son travail en 15 minutes à la fin du semestre ; les mentors accompagneront leur étudiant(e) dans la réalisation de ces tâches. %Il y aura aussi une liste de références digitales et imprimées bien choisies pour faciliter ces étapes.  

    \item[\ding{118}] Dans le cas d'un travail original, l'étudiant(e) mentoré(e) pourra \^{e}tre accompagné(e) dans la rédaction d'un article pouvant être soumis \`a un journal étudiant tel que \href{https://studc.math.ca/notes-from-the-margin/}{Notes from the Margin}, \href{https://gradmath.org/}{The Graduate Journal of Mathematics}, \href{https://msp.org/involve/about/journal/about.html}{Involve} ou \href{https://journals.calstate.edu/pump/about}{The PUMP Journal of Undergraduate Research}.

    \item[\ding{118}] Une présentation dans le cadre d'une conférence étudiante, telle que le Colloque panquébécois de l'ISM ou le Congrès canadien des étudiants en mathématiques (CC\'EM), sera encouragée.

    \item[\ding{118}] Créer des ponts et des réseaux de communication entre les étudiant(e)s intéressé(e)s à poursuivre en recherche et les professeur(e)s de l'ISM, dans le but de tirer parti des bourses de recherche estivales de l'ISM pour renforcer la communauté mathématique montréalaise en rapprochant le corps professoral des étudiant(e)s.

\end{itemize}

Le calendrier suivant donne une vue d'ensemble des activités prévues pour le projet de lecture dirigée. Il est à noter que certaines dates peuvent être sujettes à des ajustements en fonction des disponibilités des participant(e)s et des intervenant(e)s.

\medskip

\renewcommand{\arraystretch}{1.4}
\noindent\begin{tabularx}{\textwidth}{|>{\centering\arraybackslash}p{0.28\textwidth}|X|}
\hline
\rowcolor{uqamblue!85}
\textcolor{white}{\textbf{Date}} & \textcolor{white}{\textbf{Activité}} \\
\hline
Octobre -- décembre 2026 & Appel \`a candidatures auprès des étudiant(e)s de premier cycle et jumelage des étudiant(e)s avec les mentors \\
\hline
20 janvier 2027 & Séance de lancement du projet de lecture dirigée, accompagnée d'un moment convivial de détente et d'échange \\
\hline
10 février 2027 & Conférence par un(e) professeur(e) du département \\
\hline
10 mars 2027 & Conférence par un(e) professeur(e) du département \\
\hline
28 avril 2027 & Présentations des étudiant(e)s de premier cycle\\
\hline
29 avril 2027 & Présentations des étudiant(e)s de premier cycle et clôture du projet \\
\hline
30 mai 2027 & Date limite pour la remise des rapports finaux des étudiant(e)s de premier cycle \\
\hline
\end{tabularx}
\renewcommand{\arraystretch}{1}